% Options for packages loaded elsewhere
\PassOptionsToPackage{unicode}{hyperref}
\PassOptionsToPackage{hyphens}{url}
\documentclass[
  10pt,
]{article}
\usepackage{xcolor}
\usepackage[margin=1in]{geometry}
\usepackage{amsmath,amssymb}
\setcounter{secnumdepth}{-\maxdimen} % remove section numbering
\usepackage{iftex}
\ifPDFTeX
  \usepackage[T1]{fontenc}
  \usepackage[utf8]{inputenc}
  \usepackage{textcomp} % provide euro and other symbols
\else % if luatex or xetex
  \usepackage{unicode-math} % this also loads fontspec
  \defaultfontfeatures{Scale=MatchLowercase}
  \defaultfontfeatures[\rmfamily]{Ligatures=TeX,Scale=1}
\fi
\usepackage{lmodern}
\ifPDFTeX\else
  % xetex/luatex font selection
\fi
% Use upquote if available, for straight quotes in verbatim environments
\IfFileExists{upquote.sty}{\usepackage{upquote}}{}
\IfFileExists{microtype.sty}{% use microtype if available
  \usepackage[]{microtype}
  \UseMicrotypeSet[protrusion]{basicmath} % disable protrusion for tt fonts
}{}
\makeatletter
\@ifundefined{KOMAClassName}{% if non-KOMA class
  \IfFileExists{parskip.sty}{%
    \usepackage{parskip}
  }{% else
    \setlength{\parindent}{0pt}
    \setlength{\parskip}{6pt plus 2pt minus 1pt}}
}{% if KOMA class
  \KOMAoptions{parskip=half}}
\makeatother
\usepackage{longtable,booktabs,array}
\usepackage{calc} % for calculating minipage widths
% Correct order of tables after \paragraph or \subparagraph
\usepackage{etoolbox}
\makeatletter
\patchcmd\longtable{\par}{\if@noskipsec\mbox{}\fi\par}{}{}
\makeatother
% Allow footnotes in longtable head/foot
\IfFileExists{footnotehyper.sty}{\usepackage{footnotehyper}}{\usepackage{footnote}}
\makesavenoteenv{longtable}
\setlength{\emergencystretch}{3em} % prevent overfull lines
\providecommand{\tightlist}{%
  \setlength{\itemsep}{0pt}\setlength{\parskip}{0pt}}
\usepackage{bookmark}
\IfFileExists{xurl.sty}{\usepackage{xurl}}{} % add URL line breaks if available
\urlstyle{same}
\hypersetup{
  pdftitle={Specialized Complexity of Adaptive Observation},
  pdfauthor={Devin J. Wright · NIGHTFALL Research Program},
  pdfkeywords={adaptive query complexity; product probability spaces;
fractional certificates; Efron--Stein influence; Hilbert prediction;
finite-field phase retrieval; formal verification},
  hidelinks,
  pdfcreator={LaTeX via pandoc}}

\title{Specialized Complexity of Adaptive Observation}
\usepackage{etoolbox}
\makeatletter
\providecommand{\subtitle}[1]{% add subtitle to \maketitle
  \apptocmd{\@title}{\par {\large #1 \par}}{}{}
}
\makeatother
\subtitle{Certificate Geometry, Rank Structure, and Finite-Field Phase
Codes}
\author{Devin J. Wright · NIGHTFALL Research Program}
\date{11 November 2026 · research manuscript v1.0}

\begin{document}
\maketitle
\begin{abstract}
General low-interaction observability does not classify the query cost
of systems with additional structure. This paper develops a model-aware
synthesis of three complementary refinements. First, bounded fractional
certificate cost yields a coordinate sampler with an influence gain and
a geometric posterior contraction, giving an exponential squared-error
bound for Boolean observables of bounded Hoeffding degree. Second,
disjoint output-input groups permit rank-linear allocation of scalar
adaptive predictors, in contrast to the general logarithmic spectral
regime. Third, cubic-power measurements over the finite extension F16/F4
permit explicit phase coding: three fixed measurements separate four
relative phases, while a closed-branch adaptive two-measurement
construction separates three, producing sharp worst-case query counts at
dimensions four, five, and seven under the specified oracle. An
unconditional quadratic influence improvement is discussed separately as
an internal analytic candidate. Every statement specifies its sensing
model and evidentiary status; formal compilation of a lemma is not
treated as a certificate of originality or of a larger theorem.
\end{abstract}

\section{Introduction}\label{introduction}

Two general observability results form the background of this work.
\emph{Exact Recovery of Finite-Range Functions on Arbitrary Product
Domains} (Paper A) gives dimension-free, almost-sure adaptive
identification of finite-valued degree-\(d\) observables, including a
\(d^3\) whole-coordinate bound for Boolean outputs {[}A{]}.
\emph{Spectral Adaptive Prediction on Heterogeneous Product Spaces}
(Paper B) gives an author-side analytic dimension-free prediction bound
for bounded Hilbert-valued functions of any finite interaction degree,
with output covariance tail \(\tau_r\) and error tolerance \(\eta\)
{[}B{]}. Both papers go beyond the original uniform Boolean-cube
setting; their full arbitrary-measurable-product claims are analytic
manuscripts awaiting non-author proof review. The finite categorical
scalar influence ingredient used by Paper B has a separate pinned Lean
compilation {[}N1{]}.

These general results do \textbf{not} determine the best query
complexity when a system also has bounded certificates, disjoint output
channels, quadratic interactions, or a finite measurement alphabet with
adaptive phase codes. Those features make different structural resources
available. The purpose here is to identify the resulting bounds without
confusing different oracle models or silently promoting conditional
proofs.

The central organizing principle is \textbf{structure-conditioned
observation cost}. Conditional progress is measured in the correct
potential: fractional certificate overlap controls a Boolean resampling
potential, mutually independent output groups permit componentwise risk
allocation, and a finite-field powered oracle is governed by exact
transcript capacity. These are not competing proofs of a single
unqualified universal theorem.

Our source boundary is November 11, 2026. A closed research branch
containing a stronger adaptive F16 code is included with its original
October 30 attribution, but its Lean source is not part of the merged
main-branch inventory. The quadratic degree-two exponent-two claim is
treated as a separate proposed strengthening, not an input to any
unconditional implication of the other sections.

\section{Models, normalizations, and claim
levels}\label{models-normalizations-and-claim-levels}

\subsection{Whole-coordinate product
observation}\label{whole-coordinate-product-observation}

Let \(X=(X_i)_{i\in I}\) have independent coordinates on finite product
spaces, with marginal laws \(\mu_i\). A \textbf{whole-coordinate query}
reveals \(X_i\) exactly. For \(g:\prod_i\Omega_i\to\{0,1\}\), write

\[
V(g)=\operatorname{Var}(g(X)),\qquad
I_i(g)=\mathbb E\big[\operatorname{Var}(g(X)\mid X_{-i})\big]. \tag{1}
\]

Equivalently \(I_i(g)=\Pr\{g(X)=1,\,g(X^{i\leftarrow Y_i})=0\}\), where
\(Y_i\) is an independent marginal resample. The probability of
\emph{either direction} of disagreement is \(2I_i(g)\); keeping this
normalization matters. With the Hoeffding--Efron--Stein expansion
\(g=\sum_S g_S\), put \(e_S=\mathbb E g_S^2\) and

\[
T(g)=\sum_i I_i(g)=\sum_{S\ne\varnothing}|S|e_S,
\qquad V(g)\le T(g)\le dV(g) \tag{2}
\]

when interaction degree is at most \(d\). Prediction means the
conditional mean after the observed transcript, and squared-error risk
is its expected conditional variance.

For a bounded Hilbert-valued observable \(F\), replace absolute values
by Hilbert norms and define the output covariance tail
\(\tau_r=\operatorname{tr}(C_F)-\sum_{j=1}^{r}\lambda_j(C_F)\).
Whole-coordinate query counts \textbf{do not} measure computation time
or bits communicated.

\subsection{Powered finite-field
observation}\label{powered-finite-field-observation}

A separate model fixes \(K=\mathbb F_{16}=\mathbb F_2[z]/(z^4+z+1)\) and
\(H=\mathbb F_4^\times=\{1,6,7\}\) under the source's four-bit encoding.
One query selects \(a\in K^k\) and observes

\[
\mathcal O_a(x)=(a\cdot x)^3\in K. \tag{3}
\]

An adaptive coefficient vector may depend only on previous
\textbf{powered} outputs. The target is \(x\) modulo one global
multiplier in \(H\). Equation (3) is \textbf{not} a coordinate-reveal
oracle and \textbf{not} a complex magnitude-squared oracle. Their query
counts cannot be compared numerically without a reduction.

\subsection{Evidentiary categories}\label{evidentiary-categories}

We use four noninterchangeable labels throughout: \textbf{kernel-checked
finite statement}, \textbf{written analytic derivation from stated
inputs}, \textbf{exact finite computational certificate}, and
\textbf{proposed/unrefereed theorem}. A successful Lean build proves the
exact declaration under its imported hypotheses; it does not establish
historical originality, a measure-theoretic limit, or an unstated
multistep algorithm. The finite F16 certificates, the finite categorical
influence theorem, and scoped product-posterior sublemmas have
identifiable CI receipts. External mathematical refereeing and worldwide
priority clearance are not asserted for the combined results.

\section{Certificate overlap yields exponential risk
reduction}\label{certificate-overlap-yields-exponential-risk-reduction}

\subsection{Fractional rows and the overlap
invariant}\label{fractional-rows-and-the-overlap-invariant}

For Boolean \(g\) on a finite product with strictly positive marginal
atoms, a \textbf{fractional certificate row} at \(x\) is a nonnegative
vector \(\lambda^x=(\lambda_i^x)\) satisfying

\[
\sum_{i:x_i\ne y_i}\lambda_i^x\ge1
\quad\text{whenever}\quad g(y)\ne g(x),
\qquad \sum_i\lambda_i^x\le K. \tag{4}
\]

Assume \(g\) is nonconstant and \(K\ge1\). Set
\(L(x)=\sum_i\lambda_i^x\), and define class masses and loads

\[
V_b=\mathbb E\!\left[\frac{\mathbf1_{\{g=b\}}}{L(X)}\right],\quad
P_i=\mathbb E\!\left[\frac{\mathbf1_{\{g=1\}}\lambda_i^X}{L(X)^2}\right],\quad
Q_i=\mathbb E\!\left[\frac{\mathbf1_{\{g=0\}}\lambda_i^X}{L(X)^2}\right]. \tag{5}
\]

Then \(\sum_iP_i=V_1\), \(\sum_iQ_i=V_0\) and \(V_1V_0\ge V(g)/K^2\).
Introduce

\[
S=\sum_i\sqrt{P_iQ_i},\qquad
\beta=\frac{S}{\sqrt{V_1V_0}}\in(0,1],\qquad
\pi_i=\frac{\sqrt{P_iQ_i}}{S}. \tag{6}
\]

The coefficient \(\beta\) is a Bhattacharyya-type overlap of the two
normalized class-load profiles. Small overlap means that the two output
classes rely on different coordinates for certification; this
\emph{strengthens} the influence guarantee, rather than imposing a cost
in alphabet size.

\textbf{Theorem 1 (certificate-overlap influence selection; finite
input).} Suppose the product-law localized two-sided certificate
inequality

\[
V_1V_0\le\sum_i\sqrt{I_i(g)P_iQ_i} \tag{7}
\]

holds for the rows in (4). Then the explicit distribution \(\pi\) in (6)
obeys

\[
\boxed{\mathbb E_{i\sim\pi}I_i(g)\ge\frac{V(g)}{K^2\beta^2}\ge\frac{V(g)}{K^2}.} \tag{8}
\]

In particular, there is a single deterministic coordinate \(j\) with
\(I_j(g)\ge V(g)/K^2\).

\emph{Proof.} Put \(q_i=\sqrt{P_iQ_i}\), so \(\sum_iq_i=S\). Weighted
Cauchy--Schwarz and (7) give

\[
(V_1V_0)^2\le\Big(\sum_iq_i\Big)\Big(\sum_iq_i I_i(g)\Big)
=S^2\,\mathbb E_{\pi}I_i(g).
\]

Divide by \(S^2=\beta^2V_1V_0\), then use \(V_1V_0\ge V(g)/K^2\). A
maximum coordinate influence dominates its \(\pi\)-average. \(\square\)

The source of inequality (7) is NIGHTFALL's finite-product adaptation of
Turturean's all-pairs certificate mechanism; the inherited inequality
and its historical antecedents must be credited {[}N2, LZ, T{]}. Theorem
1 is the explicit geometric consequence and is not claimed to be a
globally first use of weighted Cauchy--Schwarz.

\subsection{An adaptive posterior
potential}\label{an-adaptive-posterior-potential}

\textbf{Theorem 2 (specialized Boolean observation rate).} Let \(g\) be
a Boolean function of Hoeffding degree at most \(d\ge1\) on a
\textbf{finite, full-support independent product}. Suppose all points
have valid fractional certificate rows of total cost at most \(K\ge1\),
as in (4). Using the finite localized inequality (7), a deterministic
adaptive whole-coordinate policy can achieve

\[
\boxed{\operatorname{MSE}(t)\le dV(g)\exp\!\left(-\frac{t}{dK^2}\right).} \tag{9}
\]

Consequently, for \(0<\eta<dV(g)\), a sufficient number of observations
is

\[
\boxed{N_{\mathrm{cert}}(\eta)=\left\lceil dK^2\log\frac{dV(g)}{\eta}\right\rceil,} \tag{10}
\]

or fewer if the observer has already read every available coordinate. If
\(\eta\ge dV(g)\), zero queries already suffice for this displayed
guarantee.

\emph{Proof.} After a finite transcript, conditioning on the observed
values leaves a product law on the unobserved coordinates, and its
posterior section retains interaction degree at most \(d\). Restricting
each original fractional row to unqueried coordinates preserves (4),
since two states in the same fiber agree on every observed coordinate;
its cost remains at most \(K\). At any unresolved posterior state choose
a remaining coordinate with influence at least \(V_{\mathrm{post}}/K^2\)
by Theorem 1, using a fixed-order rule for ties.

For an unqueried \(j\) the exact posterior ANOVA identity is

\[
\mathbb E_{X_j}\big[T(g\mid X_j)\big]=T(g)-I_j(g). \tag{11}
\]

Indeed, each Hoeffding component indexed by \(S\ni j\) retains its
energy in expectation but loses one unit of the cardinality weight
\(|S|\) after conditioning; cross terms vanish after the average. Since
\(V_{\mathrm{post}}\ge T_{\mathrm{post}}/d\), the chosen coordinate
gives

\[
\mathbb E[T_{t+1}\mid\mathcal T_t]\le
\left(1-\frac{1}{dK^2}\right)T_t.
\]

Iteration yields
\(\mathbb ET_t\le dV(g)(1-1/(dK^2))^t\le dV(g)e^{-t/(dK^2)}\). The
conditional-mean squared error is
\(\mathbb E V_{\mathrm{post}}\le\mathbb ET_t\). \(\square\)

This derivation is written at the level of the complete finite posterior
argument, with named Lean checks for the localized finite inequality and
an actual one-step product posterior decrement. The \textbf{entire
iterated policy is not a single Lean theorem} {[}N2{]}. The cost \(K\)
is an independent structural hypothesis; the paper does not assume
\(K=O(d^2)\) for all Boolean functions or unrestricted output alphabets.

\textbf{Corollary 2.1 (competing exact and approximate policies).} On
the Boolean finite-range subclass where Paper A's analytic \(d^3\)
exact-recovery theorem applies, one may choose the cheaper of the
exact-recovery policy and Theorem 2:

\[
N_{\mathrm{Boolean}}(\eta)\le
\min\left\{d^3,\ \left\lceil dK^2\log_+\!\frac{dV(g)}{\eta}\right\rceil\right\}, \tag{12}
\]

where \(\log_+u=\max\{0,\log u\}\). The first policy has zero error,
while the second can be preferable at coarse tolerance. This minimum is
a selection between established \emph{author-side analytic
constructions}, not a new independent optimality theorem.

\section{Disjoint output groups and linear-rank
allocation}\label{disjoint-output-groups-and-linear-rank-allocation}

The generic Hilbert spectral bound in Paper B pays roughly
\(r^{3/2}\eta^{-3/2}\) times a logarithm for arbitrary overlapping
output components. Additional independence permits a stronger allocation
principle without the unresolved quadratic theorem.

Let \(F=(g_1(X_{J_1}),\ldots,g_r(X_{J_r}))\) with pairwise disjoint
input groups \(J_j\), independent product law and degree at most \(d\).
Assume \(\|F\|\le B\) almost surely and
\(b_j=\operatorname{ess\,sup}|g_j|\). A simple near-supremum
independence argument shows

\[
\sum_{j=1}^r b_j^2\le B^2. \tag{13}
\]

Let \(c_*=\sqrt{3/12500}/\pi^6\) be the precise coefficient of the
pinned finite-product scalar influence theorem {[}N1{]}, and suppose the
arbitrary-product scalar observation corollary of Paper B is accepted at
its stated analytic level:

\[
N_{\mathrm{scal}}(b,\eta')\le
\left\lceil \frac{5d^7b^3}{3c_*(\eta')^{3/2}}\right\rceil. \tag{14}
\]

\textbf{Theorem 3 (rank-linear specialized allocation; analytic from
(14)).} Under these assumptions there exists an adaptive policy with
risk at most \(\eta\) and at most

\[
\boxed{N_{\mathrm{disjoint}}(\eta)\le
r+\frac{5d^7B^3}{3c_*\eta^{3/2}}\,r.} \tag{15}
\]

\emph{Proof.} Omit zero components. Set \(S=\sum_j b_j^{6/5}\) and
assign target risk \(\eta_j=\eta b_j^{6/5}/S\). Run an independent
scalar policy for each disjoint group. Ignoring integer ceilings, the
total is

\[
\frac{5d^7}{3c_*}\sum_j\frac{b_j^3}{\eta_j^{3/2}}
=\frac{5d^7}{3c_*\eta^{3/2}}S^{5/2}.
\]

Hölder's inequality gives \(S^{5/2}\le r(\sum_jb_j^2)^{3/2}\le rB^3\).
The ceilings contribute at most \(r\) and the errors sum to \(\eta\).
\(\square\)

This bound is \emph{linear in rank} because inputs are partitioned into
genuinely disjoint output groups. Orthogonal covariance eigenvectors by
themselves do not imply such a partition. A fixed orthogonal residual
output contributes its covariance trace \(\tau_r\) to the prediction
error. Theorem 3 is not an independent general-rank improvement over
Paper B: it records the already available specialized regime as an
explicit component of the complexity classification.

\section{The quadratic exponent-two frontier: a conditional
acceleration}\label{the-quadratic-exponent-two-frontier-a-conditional-acceleration}

The more ambitious degree-two source manuscript proposes, for
\(|f|\le B\) and Efron--Stein degree at most two under arbitrary
independent input laws,

\[
\boxed{\operatorname{Var}(f)\le C_2B\sqrt{\sup_i I_i(f)},
\quad C_2=16(1+\sqrt3).} \tag{16}
\]

This is \textbf{not} a consequence of the currently pinned
general-degree \((5/2,6)\) Lean theorem. The later conditional algebraic
closure has coefficient \(15(1+\sqrt3)\) under additional unproved
matching and fourth-moment hypotheses; it must not be substituted for
\(C_2\) in an unconditional theorem. The complete product-to-matching
construction remains an internal analytic candidate awaiting independent
specialist review and an end-to-end formal certificate {[}N3{]}.

\textbf{Proposition 4 (what exponent two would buy).} \emph{Assume} (16)
holds uniformly for every posterior section of a bounded degree-two
scalar observable, with the indicated influence normalization. A greedy
largest-influence coordinate observer satisfies

\[
\mathbb E V_{t+1}\le\mathbb EV_t
-\frac{(\mathbb EV_t)^2}{C_2^2B^2},
\qquad
\boxed{\mathbb EV_t\le\frac{C_2^2B^2}{t+C_2^2B^2/V_0}}. \tag{17}
\]

\emph{Proof.} By (16), at any posterior state some unqueried coordinate
has influence at least \(V_{\mathrm{post}}^2/(C_2^2B^2)\). The expected
reduction of conditional variance upon revealing \(j\) equals its
influence. Taking expectations and applying Jensen gives the first
inequality; reciprocating the standard quadratic-decay recursion gives
the second. \(\square\)

Thus a \emph{verified} exponent-two inequality would improve scalar
\(\eta^{-3/2}\) dependence to \(\eta^{-1}\), and the principal quadratic
manuscript further proposes a rank-\(r\) bound of order \(r/\eta\)
through a separate vector lifting argument and an additive lower-bound
family {[}N3{]}. That rank-sensitive lifting does \textbf{not} follow by
simply replacing the scalar coefficient in Paper B, and its
\(\Theta(r/\eta)\) classification is retained as an \textbf{unrefereed
candidate}, not promoted into Theorem 3.

\section{Exact finite-field phase coding and adaptive
capacity}\label{exact-finite-field-phase-coding-and-adaptive-capacity}

This section concerns the distinct oracle (3), not the
product-coordinate model. Its strength is that the key claims reduce to
finite arithmetic plus a standard global-phase normalization argument.

\subsection{Transcript capacity}\label{transcript-capacity}

The cube map on \(K\) has exactly

\[M=1+\frac{16-1}{3}=6\]

outputs. The number of vectors in \(K^k\) modulo a global \(H\)
multiplier is

\[S_k=1+\frac{16^k-1}{3}. \tag{18}\]

Any adaptive decision tree of worst-case depth \(m\) has at most \(6^m\)
transcripts. Therefore

\[
\boxed{m\ge\left\lceil\log_6 S_k\right\rceil.} \tag{19}
\]

The argument counts \emph{powered} outputs only; an unpowered
intermediate sum cannot be used to choose a future query.

\subsection{Three fixed rows for four relative
phases}\label{three-fixed-rows-for-four-relative-phases}

The merged original finite Lean certificate uses the following
\(3\times5\) coefficient matrix, with the first column an anchor:

\[
A=\begin{pmatrix}
14&12&15&3&7\\
15&4&9&6&8\\
5&5&3&14&6
\end{pmatrix}. \tag{20}
\]

For \(\eta\in H^4\) the map
\(\eta\mapsto((A(1,\eta))_1^3,(A(1,\eta))_2^3,(A(1,\eta))_3^3)\) is
injective on all \(3^4=81\) relative-phase states. This exact 81-state
statement is compiled in the merged \texttt{PhaseBlock.lean} source and
separately checked with a different finite-field multiplication
implementation {[}N4{]}. Two queries could distinguish at most
\(6^2=36\) states, so three are optimal \textbf{for this anchored
four-phase subtask}.

For a signal with \(s\ge1\) nonzero coordinates, first take all \(k\)
coordinate cubes, choose canonical roots and an anchor, then group the
\(s-1\) remaining unknown relative phases into blocks of four. Three
fixed mixed rows recover each full block; at most one query per
remainder phase suffices. The resulting written all-signal cost is

\[
N_{\mathrm{fixed}}(k,s)=k+3\left\lfloor\frac{s-1}{4}\right\rfloor+((s-1)\bmod4). \tag{21}
\]

This general transfer is recorded analytically in the original October
29 source; its finite core, not the entire decoder, is Lean-compiled.

\subsection{The stronger recovered adaptive
block}\label{the-stronger-recovered-adaptive-block}

A \textbf{pre-existing October 30 closed research branch} contains a
different, stronger adaptive gadget {[}N5{]}. The first anchored
coefficient row is \((1,2,3,4)\). For each \(\eta\in H^3\), observe
\(y=(1+2\eta_1+3\eta_2+4\eta_3)^3\) and choose a second row according to
the following \emph{measured-output} table:

\begin{longtable}[]{@{}ccr@{}}
\toprule\noalign{}
First cube \(y\) & Second coefficient row & States in first fiber \\
\midrule\noalign{}
\endhead
\bottomrule\noalign{}
\endlastfoot
\(0\) & \((0,0,0,1)\) & 1 \\
\(1\) & \((1,0,2,8)\) & 5 \\
\(8\) & \((0,1,4,8)\) & 5 \\
\(10\) & \((0,1,2,4)\) & 6 \\
\(12\) & \((0,1,8,4)\) & 5 \\
\(15\) & \((0,1,2,4)\) & 5 \\
\end{longtable}

The \(27\) ordered pairs of first and second powered outputs are
distinct. This finite Boolean certificate appears in an
\textbf{unmerged, closed} Lean source at PR \#252, with a passing
scope-specific CI receipt; the original construction and its exact
Python checker are in PR \#255. No claim is made that either branch is
merged to the main source inventory {[}N5{]}.

\textbf{Theorem 5 (adaptive cubic phase-code transfer; written algebraic
proof from the finite table).} Accept the 27-state injectivity
certificate in the table above. Every \(x\in K^k\) with support size
\(s\ge1\) can be reconstructed up to a single global phase using at most

\[
\boxed{N_{\mathrm{adapt}}(k,s)=
k+2\left\lfloor\frac{s-1}{3}\right\rfloor+((s-1)\bmod3)} \tag{22}
\]

powered measurements. The zero signal is recognized by its \(k\)
coordinate cubes. The adaptive protocol requires only legally observed
cube-power outputs.

\emph{Proof.} Obtain \(y_j=x_j^3\) for every coordinate. Since the
kernel of cubing on \(K^\times\) is \(H\), choose a canonical root
\(u_j\) of each nonzero \(y_j\), with \(x_j=\lambda_j u_j\) for an
unknown \(\lambda_j\in H\). Choose a supported anchor \(t\) and put
\(\eta_j=\lambda_j/\lambda_t\). For every triple of other supported
indices \((j_1,j_2,j_3)\), evaluate the row \((a_0,a_1,a_2,a_3)\) using
physical coefficients \(a_0/u_t,a_1/u_{j_1},a_2/u_{j_2},a_3/u_{j_3}\) on
those signal coordinates, zero elsewhere. Its unpowered sum is
\(\lambda_t(a_0+\sum_{h=1}^3a_h\eta_{j_h})\). Since \(\lambda_t^3=1\),
cubing exactly returns the table's anchored observation, independent of
the unknown anchor phase. The first powered observation selects the
second row; the second powered observation completes the unique
transcript and recovers the triple of relative phases. A leftover phase
is recovered by the single mixed form \(1+z\eta\), whose cubes are
distinct as \(\eta\) varies in \(H\): equality of cubes gives equality
up to a multiplier \(h\in H\), and linear independence of \(1,z\) over
\(\mathbb F_4\) forces \(h=1\) and \(\eta=\eta'\). Recover all
coordinates up to the common factor \(\lambda_t\). Each complete triple
uses two new queries, each remainder coordinate one. \(\square\)

The proof above is a universal mathematical \emph{deduction from a
finite certificate}. The archived Lean file does \textbf{not} compile
this entire support-selection and observation transcript construction as
one theorem.

\textbf{Corollary 5.1 (exact worst-case powered-query counts).} Under
this exact oracle and the verified finite table, the minimum worst-case
numbers of adaptive queries are

\[
\boxed{D(4)=6,\qquad D(5)=8,\qquad D(7)=11.} \tag{23}
\]

\emph{Proof.} The upper bounds are (22), maximized at full support. The
lower bounds are (19), since

\[
6^5=7776<S_4=21846,\quad
6^7=279936<S_5=349526,\quad
6^{10}=60466176<S_7=89478486.
\]

Thus no decision tree uses fewer than \(6,8,11\) queries, respectively.
\(\square\)

A distinct finite computational audit enumerates all \(4,369\)
normalized nonzero fixed four-coefficient rows and shows that every pair
of fixed rows has at least four colliding pairs among the \(27\)
anchored three-phase states; four is attained. Therefore two
\textbf{fixed} powered queries cannot do what the two \textbf{adaptive}
queries in the table do. This fixed-row noninjectivity is a
source-specific exhaustive arithmetic result supported by independent
finite-field checker implementations, \textbf{not} a claim of a global
impossibility theorem for arbitrary fixed larger codes.

\section{A model-indexed complexity
hierarchy}\label{a-model-indexed-complexity-hierarchy}

The preceding results are best read as a table of \emph{different
promises and different objectives}, not as one rate claimed for all
systems.

\begin{longtable}[]{@{}
  >{\raggedright\arraybackslash}p{(\linewidth - 4\tabcolsep) * \real{0.3333}}
  >{\raggedright\arraybackslash}p{(\linewidth - 4\tabcolsep) * \real{0.3333}}
  >{\raggedright\arraybackslash}p{(\linewidth - 4\tabcolsep) * \real{0.3333}}@{}}
\toprule\noalign{}
\begin{minipage}[b]{\linewidth}\raggedright
Observation class and structural promise
\end{minipage} & \begin{minipage}[b]{\linewidth}\raggedright
Query guarantee
\end{minipage} & \begin{minipage}[b]{\linewidth}\raggedright
Evidence scope
\end{minipage} \\
\midrule\noalign{}
\endhead
\bottomrule\noalign{}
\endlastfoot
Finite-palette degree-\(d\) whole-coordinate, arbitrary independent
product & Exact: \(d\lfloor d^2\min\{(q-1)^2,2R/\Delta\}\rfloor\) &
Paper A analytic, externally unrefereed \\
Boolean finite product, degree \(d\), uniform fractional certificate
cost \(K\) & MSE \(\le dV_0e^{-t/(dK^2)}\) & Finite inequality +
posterior composition; full policy analytic \\
General bounded rank-\(r\) Hilbert whole-coordinate & Risk
\(\le\tau_r+\eta\) at Paper B logarithmic budget & Paper B analytic;
finite scalar Lean basis \\
Hilbert output with disjoint independent groups &
\(r+\tfrac{5d^7B^3}{3c_*\eta^{3/2}}r\) & Derived analytic allocation
from Paper B \\
Scalar arbitrary-law quadratic, \emph{if} exponent-two influence holds &
\(O(C_2^2B^2/\eta)\) & Conditional proposal; not end-to-end Lean \\
Finite-field cube oracle, \(K=\mathbb F_{16}\), \(k=4,5,7\) & Exact
minimax \(6,8,11\) & Finite exact certificates + written all-signal
reduction \\
\end{longtable}

If two \emph{whole-coordinate} policies apply to the same function and
accuracy objective, their query bounds can be minimized. This includes
the Boolean exact-versus-certificate selection (12). Conversely, the
powered finite-field rows of (3) interrogate linear combinations and are
not whole-coordinate reads; including their query counts in a
cross-model minimum would be invalid.

The qualitative lesson is that \textbf{exponent, rank, certificate
geometry and alphabet capacity form separate complexity axes}. Degree
alone predicts neither exact real-output query complexity nor the
benefit of adaptive phase-coded measurements. The repeated
scalar-to-vector, finite-to-measurable and certificate-to-stopping steps
should be recorded once with the relevant proof scope and reused, not
re-proved as unrelated local lemmas.

\section{Certification, prior art, and
reproducibility}\label{certification-prior-art-and-reproducibility}

The statements above carry distinct evidence:

\begin{enumerate}
\def\labelenumi{\arabic{enumi}.}
\tightlist
\item
  \textbf{Finite product certificate transport.} The localized
  product-law all-pairs inequality and its actual one-step posterior
  consequence have named Lean receipts in the NIGHTFALL repository; the
  multistep stopping policy follows from (11) and a finite contraction
  argument, but has not been deposited as one end-to-end Lean theorem
  {[}N2{]}. Turturean, Lovett--Zhang, and earlier OSSS influence theory
  are essential antecedents {[}T,LZ,OSSS{]}.
\item
  \textbf{General Hilbert specialization.} The scalar product-law
  theorem uses the pinned Lean coefficient \(c_*\); the
  arbitrary-measurable scalar extension and scalar query bound (14) are
  analytic components of Paper B, not a separately checked full Hilbert
  Lean theorem {[}B,N1{]}. The rank-linear disjoint regime is therefore
  \textbf{analytic}, not more strongly certified than the ingredient it
  imports.
\item
  \textbf{Quadratic acceleration.} The source includes genuinely
  compiled Hoeffding decomposition, interpolation, and two-increment
  fourth-moment identities, but the actual global quadratic
  matching/assignment and many-increment transfer required for (16) are
  not compiled end-to-end. The numerical or arithmetic closure of those
  sublemmas does not establish the proposed exponent-two theorem
  {[}N3{]}.
\item
  \textbf{F16 exact code.} The 81-state merged finite block has Lean and
  separately coded arithmetic checks {[}N4{]}. The stronger adaptive
  27-state block has a scoped Lean certificate and independent
  arithmetic on closed PR branches, not on main; the written
  field-normalization proof above is separate {[}N5{]}. The exact model
  must be preserved in every lower-bound statement.
\end{enumerate}

\textbf{Independent audit plan.} First rerun the 81- and 27-state finite
checks using polynomial and quadratic-tower representations and mutate
every branch row and the power exponent. Next verify the fixed-row
\(4,369\)-normal-form exhaustion and the k4/k5/k7 counting inequalities.
Independently attack the Boolean posterior reduction by checking the
directed-influence factor, row validity under restriction, and
conservation of \(T\); then review the Holder allocation in (15),
including disjointness and near-supremum events. An external proof
reviewer must attack the unreduced analytic bridges and run a dated
same-oracle priority search before any worldwide originality assertion.
The cheapest decisive falsifiers are a legal branch collision among 27
states, a counterexample to (11), or a product-degree-two function
contradicting the proposed (16). Numerical failure-to-find is not a
proof.

This manuscript is a synthesis and refinement of prior NIGHTFALL
research checkpoints, \emph{not} an announcement of newly established
global research priority. The research direction and authorship are
attributed to Devin J. Wright; the source records explicitly credit
internal AI-assisted contributions (including Tessera and Ryev for the
archived adaptive code) and the external literature. The original source
SHAs and formal certificates remain the authoritative provenance for the
machine-checked claims.

\section{References and source
passports}\label{references-and-source-passports}

\textbf{{[}A{]}} D. J. Wright, \emph{Exact Recovery of Finite-Range
Functions on Arbitrary Product Domains}, NIGHTFALL Paper A, manuscript
(11 November 2026).
\href{https://github.com/d3lefts/d3lefts/blob/main/papers/exact-recovery/MANUSCRIPT.md}{Source}.

\textbf{{[}B{]}} D. J. Wright, \emph{Spectral Adaptive Prediction on
Heterogeneous Product Spaces}, NIGHTFALL Paper B, manuscript (11
November 2026).
\href{https://github.com/d3lefts/d3lefts/blob/main/papers/spectral-adaptive-prediction/MANUSCRIPT.md}{Source}.

\textbf{{[}N1{]}} NIGHTFALL, finite arbitrary-law product influence,
pinned Lean 4.33.0 source and CI receipt
\href{https://github.com/d3lefts/d3lefts/actions/runs/42929857470}{42929857470};
subsequent heterogeneous finite alphabets and Efron--Stein degree
bridge, PR \#325,
\href{https://github.com/d3lefts/d3lefts/actions/runs/42932085109}{CI
42932085109}.

\textbf{{[}N2{]}} NIGHTFALL, \emph{Certificate-guided adaptive
observation} (10 November 2026),
\href{https://github.com/d3lefts/d3lefts/blob/main/research/nightfall/formalization/product_space/CERTIFICATE_GUIDED_ADAPTIVE_RECOVERY_20261110.md}{source},
with pinned one-step posterior Lean
\href{https://github.com/d3lefts/d3lefts/actions/runs/42925787438}{CI
42925787438}.

\textbf{{[}N3{]}} NIGHTFALL, \emph{Adaptive Observability of
Low-Interaction Product Systems}, quadratic candidate
\href{https://github.com/d3lefts/d3lefts/blob/main/research/nightfall/observability-publication/PRINCIPAL_PAPER.md}{source};
\href{https://github.com/d3lefts/d3lefts/blob/main/research/nightfall/formalization/observability/QUADRATIC_DISCOVERY_PROOF_PASSPORT.md}{scoped
quadratic Lean proof passport}.

\textbf{{[}N4{]}} NIGHTFALL, \emph{Three-query phase compression in F16}
(29 October 2026),
\href{https://github.com/d3lefts/d3lefts/blob/main/research/nightfall/projects/adaptive-finite-field-decoder-2026-10-29/F16_THREE_QUERY_PHASE_BLOCK_20261029.md}{source},
\href{https://github.com/d3lefts/d3lefts/blob/main/research/nightfall/formalization/phase_block/PhaseBlock.lean}{finite
Lean source},
\href{https://github.com/d3lefts/d3lefts/actions/runs/40995299592}{CI
40995299592}.

\textbf{{[}N5{]}} Tessera/Ryev, \emph{Adaptive F16 phase compression}
(30 October 2026), archived
\href{https://github.com/d3lefts/d3lefts/pull/255}{Tessera PR \#255},
\href{https://github.com/d3lefts/d3lefts/pull/252}{Ryev PR \#252},
\href{https://github.com/d3lefts/d3lefts/actions/runs/41110673165}{finite
Lean CI 41110673165},
\href{https://github.com/d3lefts/d3lefts/actions/runs/41104093748}{independent
exact audit CI 41104093748}. These branches are closed and unmerged.

\textbf{{[}T{]}} D. Turturean, \emph{A proof of the Aaronson--Ambainis
conjecture} (2026),
\href{https://arxiv.org/abs/2611.09773}{arXiv:2611.09773}.

\textbf{{[}LZ{]}} S. Lovett and J. Zhang, \emph{Fractional Certificates
for Bounded Functions}, ITCS 2023,
\href{https://drops.dagstuhl.de/doi/10.4230/LIPIcs.ITCS.2023.84}{published
version}.

\textbf{{[}OSSS{]}} R. O'Donnell, M. Saks, O. Schramm and R. Servedio,
\emph{Every Decision Tree Has an Influential Variable}, FOCS 2005,
\href{https://www.cs.columbia.edu/~rocco/papers/focs05osss.html}{abstract}.

\textbf{{[}GJN{]}} G. Grimmett, S. Janson and J. Norris, \emph{Influence
in Product Spaces}, \emph{Advances in Applied Probability} 48 (2016),
\href{https://doi.org/10.1017/apr.2016.46}{DOI}.

\textbf{{[}CWY{]}} M. Charikar, X. Wu and Y. Ye, \emph{Adaptive Discrete
Phase Retrieval}, SOSA 2020,
\href{https://doi.org/10.1137/1.9781611976014.9}{DOI}. This uses complex
intensity measurements, not the finite cube-power oracle (3).

\end{document}
